% Generated draft: requires mathematical and source review.
\documentclass[11pt]{article}
\usepackage[margin=25mm]{geometry}
\usepackage{fontspec}
\setmainfont{Latin Modern Roman}
\usepackage{amsmath,amssymb,mathtools,unicode-math}
\setmathfont{Latin Modern Math}
\usepackage{enumitem,needspace,xurl,hyperref}
\setlength{\parindent}{0pt}
\setlength{\parskip}{5pt}
\setlength{\emergencystretch}{3em}
\newcommand{\N}{\mathbb{N}}
\newcommand{\Z}{\mathbb{Z}}
\newcommand{\Q}{\mathbb{Q}}
\newcommand{\R}{\mathbb{R}}
\newcommand{\C}{\mathbb{C}}
\newcommand{\F}{\mathbb{F}}
\newcommand{\A}{\mathbb{A}}
\newcommand{\PP}{\mathbb{P}}
\newcommand{\E}{\mathbb{E}}
\newcommand{\eps}{\varepsilon}
\newcommand{\dd}{\,\mathrm{d}}
\begin{document}
\section*{UnsolvedMath ID: 30000949}
\textbf{Source problem code:} catalogue entry
\subsection{Definitions and mathematical statement}
A prosoluble group is a profinite group expressible as an inverse limit of finite soluble groups. An abstract group $H$ is perfect if $H=[H,H]$, where $[H,H]$ is the subgroup generated by commutators, without taking a topological closure.

The substantive abstract-image interpretation of the source's question is
\[
\exists G\text{ prosoluble}\ \exists H\ne\{1\}\ \exists\pi:G\twoheadrightarrow H
\text{ an abstract homomorphism},\qquad H=[H,H].
\]
Equivalently, is there a prosoluble $G$ and a proper normal subgroup $N\lhd G$, not required to be closed, such that
\[
G=N[G,G]?
\]
The image is not required to carry a profinite quotient topology. No finite or topological finite-generation hypothesis is imposed on $G$.

Continuity matters: a continuous quotient in the category of profinite groups is again prosoluble. If it were perfect, every finite quotient would be both soluble and perfect and hence trivial, so the profinite quotient would be trivial. Thus the question cannot be interpreted as seeking a nontrivial continuous profinite perfect quotient. The extract does not explicitly say ``abstract image''; that convention is made explicit rather than treating the source as an already-impossible continuous-quotient question.
\subsection{Short English statement}
Can a prosoluble group have a nontrivial perfect abstract quotient, allowing a nonclosed kernel? The distinction from continuous profinite quotients is essential.
\subsection{Sources}
\begin{itemize}
\item Supplied report: \emph{Groups and Geometries} (2009), Oberwolfach Reports, DOI 10.4171/owr/2008/20, pp. 1031--1078; proposed year 2008. \url{https://doi.org/10.4171/owr/2008/20}.
\item Supplied reference: Nikolov, profinite group questions, \url{https://ems.press/content/serial-article-files/37105}.
\item Local UnsolvedMath registry, permanent ID 30000949, catalogue entry; created 25 August 2026, updated 16 September 2026. Its exact-statement note reports manual confirmation of coherence and absence of layout artifacts. The supplied 21 August 2026 triage labels the general existence question open and reports strong negative finitely generated cases and related classes, without specifying their full hypotheses. This dated assessment is not independently verified here.
\end{itemize}
\end{document}
